\documentclass[11pt,a4paper]{article}

% ── Packages ───────────────────────────────────────────────────────────────────
\usepackage[margin=2.2cm]{geometry}
\usepackage{amsmath,amssymb}
\usepackage{booktabs}
\usepackage{array}
\usepackage{xcolor}
\usepackage{listings}
\usepackage{tikz}
\usetikzlibrary{shapes.geometric, arrows.meta, positioning}
\usepackage{hyperref}
\usepackage{microtype}
\usepackage{enumitem}
\usepackage{fancyhdr}
\usepackage{titlesec}
% ── Colors ─────────────────────────────────────────────────────────────────────
\definecolor{ppoBlue}{RGB}{41,128,185}
\definecolor{cfOrange}{RGB}{211,84,0}
\definecolor{darkGray}{RGB}{44,62,80}
\definecolor{lightBg}{RGB}{245,247,250}
\definecolor{accentRed}{RGB}{192,57,43}
\definecolor{green}{RGB}{39,174,96}

% ── Section styling ────────────────────────────────────────────────────────────
\titleformat{\section}{\large\bfseries\color{darkGray}}{§\thesection}{0.8em}{}[{\color{ppoBlue}\titlerule[0.8pt]}]
\titleformat{\subsection}{\normalsize\bfseries\color{ppoBlue}}{\thesubsection}{0.6em}{}

% ── Header / footer ───────────────────────────────────────────────────────────
\pagestyle{fancy}
\fancyhf{}
\lhead{\small\textbf{RL Dynamic Request Batching} — Project Guide}
\rhead{\small Sudarshan Sudhakar, May 2026}
\cfoot{\thepage}
\renewcommand{\headrulewidth}{0.4pt}

% ── Code listings ──────────────────────────────────────────────────────────────
\lstset{
  basicstyle=\ttfamily\small,
  keywordstyle=\color{ppoBlue}\bfseries,
  commentstyle=\color{gray}\itshape,
  stringstyle=\color{accentRed},
  backgroundcolor=\color{lightBg},
  frame=single, framerule=0.4pt,
  rulecolor=\color{ppoBlue!40},
  breaklines=true,
  numbers=none,
  xleftmargin=6pt, xrightmargin=6pt,
}

% ── Highlight boxes ────────────────────────────────────────────────────────────
\newmdenv[
  backgroundcolor=ppoBlue!8,
  linecolor=ppoBlue,
  linewidth=1pt,
  roundcorner=4pt,
  innertopmargin=6pt, innerbottommargin=6pt,
  innerleftmargin=8pt, innerrightmargin=8pt,
]{infobox}

\newmdenv[
  backgroundcolor=cfOrange!8,
  linecolor=cfOrange,
  linewidth=1pt,
  roundcorner=4pt,
  innertopmargin=6pt, innerbottommargin=6pt,
  innerleftmargin=8pt, innerrightmargin=8pt,
]{warnbox}

% ── Hyperref ──────────────────────────────────────────────────────────────────
\hypersetup{colorlinks=true, linkcolor=ppoBlue, urlcolor=ppoBlue}

% ══════════════════════════════════════════════════════════════════════════════
\begin{document}
% ══════════════════════════════════════════════════════════════════════════════

% ── Title ──────────────────────────────────────────────────────────────────────
\begin{center}
  {\LARGE\bfseries\color{darkGray}%
   RL-Based Dynamic Request Batching\\[4pt]
   \large on GPU Inference Servers}\\[8pt]
  {\large Sudarshan Sudhakar \quad \texttt{sudarshansudhakar@gmail.com}}\\[4pt]
  {\normalsize Semester 6 Reinforcement Learning Project $\;\cdot\;$ May 2026}
  \vspace{8pt}
  \hrule height 1.4pt
\end{center}

\vspace{4pt}
\begin{infobox}
\textbf{One-sentence summary.}\; A PPO agent learns \emph{when} to dispatch
queued inference requests to a GPU backend, outperforming the production
Cloudflare dual-threshold heuristic by adapting dispatch timing dynamically
to queue depth, urgency, EMA arrival rate, and time-of-day.
\end{infobox}
\vspace{8pt}

% ══════════════════════════════════════════════════════════════════════════════
\section{Problem}
% ══════════════════════════════════════════════════════════════════════════════

A GPU inference node receives requests at \textbf{2\,000 req/s} (peak 5\,000).
Every \textbf{10\,ms} a batching layer decides to \emph{Wait} or \emph{Serve}.

\medskip
\renewcommand{\arraystretch}{1.25}
\begin{tabular}{@{}lll@{}}
\toprule
\textbf{Decision} & \textbf{Effect} & \textbf{Problem} \\
\midrule
Serve too early & tiny batch & wasted GPU kernel launch, poor utilisation \\
Wait too long   & large batch & high latency, SLA violations \\
Optimal         & adaptive    & threshold depends on rate, age, time-of-day \\
\bottomrule
\end{tabular}

\medskip
The SLA constraint is on \emph{total client-perceived latency}:
\begin{equation}
  \text{total latency} = \underbrace{w_i}_{\text{queue wait}} +
  \underbrace{t_\text{proc}(n)}_{\text{GPU time}}
  = w_i + 8 + 0.12\,n \;\text{ms}
  \;\leq\; 500\;\text{ms}
  \label{eq:sla}
\end{equation}
No closed-form optimal policy exists; fixed thresholds cannot adapt to
time-varying traffic. This motivates a learned policy.

% ══════════════════════════════════════════════════════════════════════════════
\section{MDP Formulation}
% ══════════════════════════════════════════════════════════════════════════════

\subsection{State Space — 8-dimensional continuous}

\begin{tabular}{@{}clll@{}}
\toprule
\textbf{Dim} & \textbf{Feature} & \textbf{Range} & \textbf{Role} \\
\midrule
0 & \texttt{pending\_requests} & $[0, 512]$      & Queue depth \\
1 & \texttt{oldest\_wait\_ms}  & $[0, 1000]$     & Head-of-queue age \\
2 & \texttt{request\_rate}     & $[0, 25\,000]$  & EMA arrival rate (req/s) \\
3 & \texttt{since\_serve\_ms}  & $[0, 1000]$     & Time since last dispatch \\
4 & \texttt{batch\_fill\_ratio}& $[0, 1]$        & \texttt{pending / max\_batch} \\
5 & \texttt{time\_of\_day}     & $[0, 24)$       & Fractional hour \\
6 & \texttt{urgency\_ratio}    & $[0, 3]$        & $w_\max / (\text{SLA} - t_\text{proc}(n))$ \\
7 & \texttt{rate\_trend}       & $[-1, 1]$       & Fractional EMA change per step \\
\bottomrule
\end{tabular}

\medskip
\noindent\textbf{Key design choices:}
\begin{itemize}[noitemsep]
  \item \textbf{\texttt{urgency\_ratio}}: when $>1.0$, the current batch
        \emph{will} violate SLA even if served now — gives the agent pre-emptive
        signal before violations occur.
  \item \textbf{\texttt{rate\_trend}}: detects incoming traffic spikes before
        the EMA fully adapts, enabling proactive dispatch.
  \item \textbf{EMA} smooths the instantaneous rate with $\alpha=0.1$ per 10\,ms step.
\end{itemize}

\subsection{Action Space}

$\mathcal{A} = \{0\;\text{(Wait)},\;1\;\text{(Serve)}\}$ — binary, Discrete(2).

\subsection{Reward Function}

\textbf{Serve} ($a=1$):
\begin{equation}
  r = \alpha\,n - \beta\,w_\max - c_\text{dispatch} -
  \lambda \sum_{i=1}^{n} \max\!\left(0,\; w_i + t_\text{proc}(n) - \text{SLA}\right)
\end{equation}

\textbf{Wait} ($a=0$):
\begin{equation}
  r = -\beta\,w_\max - \mathbf{1}_{n=0}\,c_\text{idle}
\end{equation}

\medskip
\begin{tabular}{@{}llll@{}}
\toprule
\textbf{Symbol} & \textbf{Value} & \textbf{Symbol} & \textbf{Value} \\
\midrule
$\alpha$ (revenue/req) & 1.0 & $c_\text{dispatch}$ (GPU overhead) & 5.0 \\
$\beta$ (urgency weight) & 0.003 & $\lambda$ (SLA penalty/ms) & 0.3 \\
$c_\text{idle}$ (empty queue) & 0.5 & & \\
\bottomrule
\end{tabular}

\medskip
\begin{warnbox}
\textbf{Critical.} $c_\text{dispatch} = 5.0$ means serving a batch of fewer than
5 requests is unprofitable ($\alpha n < c$). This single hyperparameter drives
batch accumulation and was tuned empirically.
\end{warnbox}

% ══════════════════════════════════════════════════════════════════════════════
\section{Algorithms}
% ══════════════════════════════════════════════════════════════════════════════

\subsection{PPO — Proximal Policy Optimization (Primary)}

On-policy actor-critic using the clipped surrogate objective:
\begin{equation}
  L^\text{CLIP} = \mathbb{E}_t\!\left[
    \min\!\left(r_t(\theta)\,\hat{A}_t,\;
    \text{clip}(r_t(\theta), 1{-}\epsilon, 1{+}\epsilon)\,\hat{A}_t
    \right)
  \right], \quad \epsilon = 0.2
\end{equation}
where $r_t(\theta) = \pi_\theta(a_t|s_t) / \pi_{\theta_\text{old}}(a_t|s_t)$.
Advantages are estimated via GAE ($\lambda=0.95$).

\medskip
\begin{tabular}{@{}llll@{}}
\toprule
\textbf{Hyperparameter} & \textbf{Value} & \textbf{Hyperparameter} & \textbf{Value} \\
\midrule
Learning rate & $3\times10^{-4}$ & Entropy coefficient & 0.005 \\
$n_\text{steps}$ & 2\,048 & GAE $\lambda$ & 0.95 \\
Minibatch size & 256 & Discount $\gamma$ & 0.99 \\
Gradient epochs & 10 & Architecture & MLP $[128,128]$ \\
Parallel envs & 4 & Total timesteps & 1\,000\,000 \\
\bottomrule
\end{tabular}

\subsection{D3QN — Dueling Double DQN + Prioritized Replay}

Combines three improvements over vanilla DQN:
\begin{enumerate}[noitemsep]
  \item \textbf{Dueling streams}: $Q(s,a) = V(s) + A(s,a) - \frac{1}{|\mathcal{A}|}\sum_{a'}A(s,a')$
  \item \textbf{Double DQN}: action selected by online network, evaluated by target
        network — removes overestimation bias
  \item \textbf{PER}: sample probability $p_i \propto |\delta_i|^\alpha$
        ($\alpha=0.6$), IS-corrected with $\beta$: $0.4\to1.0$
\end{enumerate}
Additionally uses \textbf{3-step returns} to propagate signal across the
wait$\to$accumulate$\to$serve cycle.

\subsection{Discrete SAC}

Maximum-entropy framework with automatic temperature tuning.
Target entropy $H^* = 0.5\ln(2) \approx 0.347$ nats.
Soft Q-function updates with $\tau=0.005$ target smoothing.
Reward scaled by $10^{-3}$ to keep Q-values in $[3,5]$.

\subsection{Recurrent PPO (Ablation)}

PPO augmented with an LSTM layer before policy and value heads.
Handles temporal correlations explicitly but requires stateful
inference and $\sim$2$\times$ more training time.
Used to verify that the MLP's EMA-based features are sufficient.

\medskip
\begin{infobox}
\textbf{Why PPO wins.}
The 8 state features already encode sufficient history (EMA rate, rate trend,
since-serve, oldest-wait). An explicit LSTM is unnecessary. PPO trains in
$\sim$25\,min on CPU, produces a stateless policy, and deploys with a
single \texttt{predict()} call.
\end{infobox}

% ══════════════════════════════════════════════════════════════════════════════
\section{Baselines}
% ══════════════════════════════════════════════════════════════════════════════

\subsection{Cloudflare Dual-Threshold (Competitive)}

The actual production heuristic used by Cloudflare AI Gateway, NVIDIA Triton,
and vLLM:
\begin{alignat}{2}
  \text{SERVE}&\;\text{ if }\; w_\max &\;\geq\;& 0.80 \times (\text{SLA} - t_\text{proc}(n))
    \quad\text{[latency urgency]}\\
  \text{SERVE}&\;\text{ if }\; n      &\;\geq\;& \lfloor\text{rate} \times 0.050\rfloor
    \quad\text{[batch efficiency]}
\end{alignat}
At 2\,000 req/s: batch target = 100. At 5\,000 req/s: target = 250.
\emph{Adaptive to load but blind to rate trajectory, time-of-day, and
inter-serve interval.}

\subsection{GreedyBatch (Floor)}

Serves whenever $n \geq 8$. Minimises latency but maximises dispatch overhead
and wastes GPU capacity — provides a lower bound on efficiency.

\begin{warnbox}
\textbf{Why the textbook $e^{-\lambda t}$ formula was discarded.}
At $\lambda=2000$ req/s and SLA = 0.5\,s: $e^{-2000 \times 0.5} \approx 0$.
The agent never dispatches, violates every SLA, and scores $-6000+$ per episode.
The formula is designed for web batching at $\lambda \approx 1$--10 req/s and
does not transfer to inference serving.
\end{warnbox}

% ══════════════════════════════════════════════════════════════════════════════
\section{Implementation Details}
% ══════════════════════════════════════════════════════════════════════════════

\subsection{Environment (\texttt{env/batching\_env.py})}

\begin{itemize}[noitemsep]
  \item \textbf{Traffic generator}: Poisson arrivals with time-of-day multipliers
        ($\times2.5$ peak hours 9--23, $\times0.2$ off-peak). Each episode
        samples a random start hour from $[0,24)$.
  \item \textbf{Queue}: Python list of arrival timestamps; capped at 512 entries.
        Overflow sheds \emph{new} arrivals (oldest requests preserved for SLA reasons).
  \item \textbf{GPU model}: $t(n) = 8.0 + 0.12n$ ms. Linear fit to NVIDIA Triton
        benchmarks on an A10G (batch=1: 8.1\,ms; batch=256: 38.7\,ms).
  \item \textbf{Metrics}: P50/P95/P99 total latency and SLA violation rate
        accumulated per episode and returned in \texttt{info}.
\end{itemize}

\subsection{Training (\texttt{agent/train.py})}

\begin{lstlisting}[language=Python]
env = make_vec_env(BatchingEnv, n_envs=4, vec_env_cls=SubprocVecEnv)
env = VecNormalize(env, norm_obs=True, norm_reward=False)

model = PPO("MlpPolicy", env,
            learning_rate=3e-4, n_steps=2048, batch_size=256,
            n_epochs=10, gamma=0.99, gae_lambda=0.95,
            clip_range=0.2, ent_coef=0.005,
            policy_kwargs={"net_arch": [128, 128]})

model.learn(total_timesteps=1_000_000,
            callback=EvalCallback(eval_env, best_model_save_path="models/best",
                                  eval_freq=20_000, n_eval_episodes=10))
model.save("models/ppo_final")
env.save("models/ppo_vecnorm.pkl")
\end{lstlisting}

\subsection{Deployment (\texttt{deploy/middleware.py})}

\begin{lstlisting}[language=Python]
from deploy.middleware import BatchingMiddleware

mw = BatchingMiddleware(
    model_path   = "models/ppo_final",
    vecnorm_path = "models/ppo_vecnorm.pkl",
)

# Called every 10 ms
if mw.should_dispatch():
    batch = mw.flush()
    inference_backend.process(batch)
\end{lstlisting}

\begin{warnbox}
\textbf{Critical}: \texttt{ppo\_final.zip} and \texttt{ppo\_vecnorm.pkl} store
running normalization statistics jointly with the policy. Always regenerate
\emph{both} files when retraining. Using stale VecNormalize stats with new
weights produces silently incorrect predictions.
\end{warnbox}

% ══════════════════════════════════════════════════════════════════════════════
\section{File Structure}
% ══════════════════════════════════════════════════════════════════════════════

\begin{lstlisting}
config.py                  Central config: env, GPU model, PPO/D3QN/SAC params
env/
  batching_env.py          Gymnasium env (8-dim state, Poisson traffic, rewards)
  traffic_generator.py     Time-of-day Poisson arrival generator
baselines/
  cloudflare_formula.py    Cloudflare dual-threshold + GreedyBatch
agent/
  train.py                 PPO training (VecNormalize, 4 envs, 1M steps)
  train_d3qn.py            D3QN training (PER, dueling, 3-step returns)
  train_sac.py             Discrete SAC training
  d3qn.py                  D3QN network architecture
  discrete_sac.py          Discrete SAC implementation
  evaluate.py              3-agent comparison, generates plots
deploy/
  middleware.py            Production class: loads model + VecNorm stats
models/
  ppo_final.zip            Trained PPO weights
  ppo_vecnorm.pkl          VecNormalize statistics (deploy alongside ppo_final)
  best/best_model.zip      Best checkpoint by eval reward
results/
  comparison.png           6-panel comparison figure
  decision_heatmap.png     PPO P(Serve) heatmap with Cloudflare overlay
tensorboard_logs/          TensorBoard training curves
\end{lstlisting}

% ══════════════════════════════════════════════════════════════════════════════
\section{Quick-Start Commands}
% ══════════════════════════════════════════════════════════════════════════════

\begin{lstlisting}[language=bash]
# Install dependencies
pip install "stable-baselines3[extra]" gymnasium scipy matplotlib numpy

# Sanity check: validate env + run baseline eval
python3 test_baselines.py

# Train PPO (~20-40 min on CPU)
python3 agent/train.py

# Monitor training (open localhost:6006)
tensorboard --logdir tensorboard_logs/

# Evaluate all agents and generate plots
python3 agent/evaluate.py
# outputs: results/comparison.png, results/decision_heatmap.png

# Evaluate using best checkpoint (not final)
python3 agent/evaluate.py --model models/best/best_model

# Smoke-test deployment middleware
python3 deploy/middleware.py
\end{lstlisting}

% ══════════════════════════════════════════════════════════════════════════════
\section{Evaluation Metrics and Interpretation}
% ══════════════════════════════════════════════════════════════════════════════

\begin{tabular}{@{}lll@{}}
\toprule
\textbf{Metric} & \textbf{Higher/Lower is better} & \textbf{What it measures} \\
\midrule
Mean episode reward & Higher & Throughput $-$ latency cost $-$ SLA penalties \\
P50 total latency & Lower & Typical client experience \\
P95/P99 total latency & Lower & Tail latency (SLA-relevant) \\
SLA violation rate & Lower & Fraction of requests $>$ 500\,ms \\
Mean batch size & Higher (to a point) & GPU utilisation \\
\bottomrule
\end{tabular}

\medskip
\noindent\textbf{What to look for:}
\begin{enumerate}[noitemsep]
  \item PPO reward should exceed Cloudflare by 10--25\% — meaningful but not
        implausible. If much higher, check for reward shaping bugs.
  \item Decision heatmap: PPO boundary should be \emph{curved} (non-linear);
        Cloudflare shows two straight threshold lines.
  \item GreedyBatch: lowest P99 latency but lowest reward due to tiny batches.
\end{enumerate}

% ══════════════════════════════════════════════════════════════════════════════
\section{Experiment Presets (Eval Only)}
% ══════════════════════════════════════════════════════════════════════════════

\begin{tabular}{@{}llll@{}}
\toprule
\textbf{Preset} & \textbf{Rate} & \textbf{Peak mult.} & \textbf{Use case} \\
\midrule
\texttt{off\_peak}    & 400 req/s   & 1.5$\times$ & Overnight, minimal load \\
\texttt{standard}     & 2\,000 req/s & 2.5$\times$ & Training distribution \\
\texttt{peak\_load}   & 5\,000 req/s & 3.0$\times$ & Business hours stress test \\
\texttt{viral\_spike} & 20\,000 req/s & —          & Burst / DDoS scenario \\
\bottomrule
\end{tabular}

% ══════════════════════════════════════════════════════════════════════════════
\section{Limitations \& Future Work}
% ══════════════════════════════════════════════════════════════════════════════

\begin{itemize}[noitemsep]
  \item \textbf{Simulated GPU model}: $t(n) = 8 + 0.12n$ is a linear fit;
        real GPU throughput is non-linear (memory-bandwidth saturation effects
        at large batches, kernel fusion at medium batches).
  \item \textbf{Single node}: does not handle multi-replica load balancing or
        heterogeneous GPU tiers.
  \item \textbf{Homogeneous requests}: ignores variable request complexity
        (e.g., different sequence lengths in LLM inference).
  \item \textbf{Online adaptation}: the deployed model does not fine-tune on
        production traffic; distribution shift may degrade performance over time.
  \item \textbf{Next steps}: real Triton profiling data for GPU model;
        multi-node coordination; heterogeneous SLA tiers; online RL from
        production feedback.
\end{itemize}

\vspace{10pt}
\hrule height 0.8pt
\vspace{4pt}
\begin{center}
  \small\textit{Project repo:}
  \texttt{/Desktop/Semester\,6/Rl\,project\,gooooo/}\quad
  \textit{Model artifacts:} \texttt{models/ppo\_final.zip} +
  \texttt{models/ppo\_vecnorm.pkl}
\end{center}

\end{document}
